where the meaning of the arrows is evident, and whose commutativity, by 1.10, need only be verified on the elements @@M0@@ of @@M1@@ (resp. @@M2@@ of @@M3@@), a verification that is immediate.

By 1.8, the product of two series under this law is calculated by substituting their coefficients in the universal polynomials of the law. We denote the @@M4@@-th universal polynomial by @@M5@@; it is a polynomial in @@M6@@, isobaric of weight @@M7@@ in each of the two groups of variables (see 1.15), and symmetric with respect to these two groups. As an example, let us give the value of @@M8@@:
@@M9@@

We also have the following relation: let @@M10@@, @@M11@@, @@M12@@, then
\sgaStableTarget{sga6.e05.eq.2-3-2}@@M13@@
As usual, the standard functorial argument reduces it to the case @@M14@@, @@M15@@, @@M16@@, where it is precisely 2.3.1.

Applying 2.3.2 to @@M17@@, and to @@M18@@, we see that the element @@M19@@ is the unit for this multiplication. To avoid confusion, we sometimes denote it by @@M20@@, or by @@M21@@.

The functor @@M22@@ can now be regarded as a functor from @@M23@@ to the category of commutative rings with unit, and can therefore be composed with itself; to avoid confusion,
we
